\chapter{Fase 5: el Pronóstico, ¿cuándo conviene ir?}
\label{cap:d-f5}
\textbf{Fechas}: 30 de septiembre y 1 de octubre de 2026. \textbf{Notas}: \codigo{11-pronostico.md}, \codigo{12-planeador.md},
\codigo{15-norte-en-planeador.md}.

\textbf{Qué se buscaba}: cuántos visitantes se esperan en cada lugar los próximos 12 meses, con un rango; la probabilidad de
tormenta; escenarios malo, probable y bueno; y el riesgo de rebasar lo que cada lugar puede recibir.

\begin{opciones}[5.1 · Qué se pronostica: ``Medidas + norte'' · 30-sep-2026]
\begin{center}
\small
\begin{tabularx}{\linewidth}{>{\raggedright\arraybackslash}p{0.3\linewidth}Y}
\encabezado \h{Opción} & \h{Por qué sí o por qué no} \\
\textbf{Medidas + norte (elegida)} & Visitantes INAH de la Bahía y de la Ruta (127 meses), cruces de Belice por Chetumal (90)
y la ocupación de Cancún como referencia (55): solo series medidas que llegan a 2026 \\
\filagris Solo INAH & Chetumal se quedaría sin serie y no se sabría cuándo se llena el norte \\
Estimar la ocupación del sur & Contradice la decisión 3.4 y solo hay 36 meses de Chetumal \\
\bottomrule
\end{tabularx}
\normalsize
\end{center}
\textbf{Huecos declarados}: Maya Ka'an y la Laguna no tienen serie; el Tren Maya tiene solo 20–23 meses (para aprender la
temporada hacen falta 24), así que se usa como variable, no se pronostica.
\end{opciones}

\begin{opciones}[5.2 · Los meses que no deben enseñar: ``Hueco + forma del año'' · 30-sep-2026]
\textbf{La pregunta}: un mes cerrado no es ``cero demanda''. Si el modelo aprende de esos ceros, creerá que en abril puede no
llegar nadie por temporada.

\textbf{Elegida}: los meses de cierre, los meses parciales (justo antes de cerrar o después de reabrir) y la pandemia se
marcan y \textbf{no entrenan}, pero su valor se conserva. La forma del año se aprende de todos los años abiertos y el nivel,
de los meses posteriores a la reapertura.

\textbf{Descartados}: entrenar solo desde la reapertura (19 meses, menos de los 24 necesarios) y tomar los ceros como dato
(enseña una temporada falsa).

\textbf{Evidencia}: Kohunlich tuvo 21,850 visitantes en 2025 contra 42,813 en 2019: reabrió a la mitad. Una región toma el
motivo de su zona más afectada (en enero de 2025 Kohunlich ya abría pero Dzibanché no: sumar solo las abiertas bajaría el
total de forma artificial).
\end{opciones}

\begin{opciones}[5.3 · La forma del año: clásica y no STL · 30-sep-2026]
\textbf{Decisión técnica}: descomposición clásica multiplicativa sobre años completos. El plan decía STL, pero STL necesita
una serie continua y tiraría los años posteriores a los huecos. STL quedó como segunda opinión (coincide entre 0.85 y
0.97). Es multiplicativa porque la temporada escala con el nivel.

\textbf{Resultado}: la Ruta recibe en enero 61\,\% más que un mes promedio (1.606) y en septiembre la mitad (0.496). Las zonas
del sur se mueven con Cancún (parecido 0.83 y 0.72); Belice, no (0.26).
\end{opciones}

\begin{opciones}[5.4 · Hasta cuándo dura la pandemia en Belice · 30-sep-2026]
\textbf{Lo que pasó}: con la pandemia terminando en febrero de 2022, STL confundía la recuperación de la frontera con
temporada (ponía marzo en 0.73, cuando los años completos lo ponen entre 0.92 y 1.18; parecido de solo 0.446). De marzo a
junio de 2022 los cruces iban en 65.1, 69.5, 80.2 y 82.2\,\% de 2019; desde julio, nunca bajaron de 86\,\%.

\textbf{Elegida por Brandon}: la pandemia de Belice termina en \textbf{junio de 2022}. El parecido sube a \textbf{0.847}.
Belice entrena 62 meses en lugar de 66.

\textbf{Precisión que se dejó escrita}: la opción se presentó como ``cuando los cruces vuelven a 80\,\% o más''; mayo ya
estaba en 80.2\,\%, así que por ese criterio literal el corte sería abril. El de junio se sostiene en el salto a 86\,\% o
más desde julio.
\end{opciones}

\begin{opciones}[5.5 · Cómo elegir el modelo de pronóstico · 30-sep-2026]
\textbf{Contexto}: se compararon 5 modelos ``viajando al pasado'' (origen móvil): la línea base (el mismo mes del año
anterior), dos Holt-Winters con la forma del año fija (con y sin tendencia), una regresión con clima y Gradient Boosting.
Cada uno con un rango del 90\,\% conformal y su cobertura medida.

\begin{center}
\small
\begin{tabularx}{\linewidth}{>{\raggedright\arraybackslash}p{0.3\linewidth}Y}
\encabezado \h{Opción} & \h{Por qué sí o por qué no} \\
\textbf{Menor error con rango ≥ 80\,\% (elegida)} & Un rango ``del 90\,\%'' que falla mucho engaña a quien planea \\
\filagris Menor error sin importar el rango & En la Ruta elegiría un modelo cuyo rango acierta solo 72 de cada 100 veces \\
Un solo modelo para todo & En Cancún la regresión se equivoca 58\,\% más que la línea base \\
\bottomrule
\end{tabularx}
\normalsize
\end{center}

\textbf{Resultado}: regresión con clima en la Bahía (error 15.7\,\%, rango cumplido 91.3\,\%), la Ruta (21.6\,\%; 80.3\,\%) y
Belice (11.0\,\%; 94.1\,\%); línea base en Cancún (3.8\,\%; 89.3\,\%). Ejemplo: Bahía, diciembre de 2026, 1,311 visitantes
(entre 959 y 1,793).

\textbf{Decisiones técnicas que se declararon}: la forma del año se fija porque desde las reaperturas solo hay 17 a 20 meses
seguidos; la variante sin tendencia se agregó después de ver que la tendencia pronosticaba $-240$ visitantes. La Ruta no
llega a 90\,\% con ningún modelo (en 2023 cayó 10.7\,\% sin aviso), y el clima casi no mejora el pronóstico.
\end{opciones}

\begin{opciones}[5.6 · El golpe de una tormenta: ``Supuesto con barrido'' · 30-sep-2026]
\textbf{La pregunta}: la probabilidad de tormenta sí está medida (31 eventos, Poisson: agosto 13.9\,\%, al menos una en el año
40.3\,\%). Pero ¿cuántos visitantes quita? Solo hubo 3 meses con tormenta entre los que entrenan (Earl 2016, Franklin 2017,
Lisa 2022).

\textbf{Elegida}: un supuesto con tres valores: 0\,\%, 25\,\% y 50\,\%. \textbf{Descartado}: solo mostrar la probabilidad, sin
golpe en el escenario malo.

\textbf{Resultado}: con el golpe más duro, el escenario probable del año baja solo 1.2–2.3\,\%: las tormentas pesan en el mes,
no en el año.
\end{opciones}

\begin{opciones}[5.7 · ¿La campaña entra en los escenarios? · 30-sep-2026]
\textbf{Elegida}: \textbf{no, se ve en la Fase 6}. Los escenarios son ``lo que pasaría de todos modos''. \textbf{Descartado}:
meter un efecto supuesto de +5\,\% o +10\,\%, que mezclaría un pronóstico medido con un número sin fuente.
\end{opciones}

\begin{opciones}[5.8 · La capacidad de cada lugar: ``Capacidad probada'' · 30-sep-2026]
\textbf{Elegida}: el mes más alto que cada lugar ya recibió: Ruta 10,465 (enero de 2018), Bahía 1,639 (abril de 2017) y Belice
65,792 (agosto de 2025).

\textbf{Descartados}: cuartos de hotel (se pronostican visitantes a zonas y cruces, no noches; habría que suponer cuántos se
hospedan) y dejarlo para la Fase 6.

\textbf{Resultado del Monte Carlo} (10,000 futuros): riesgo de rebasar la capacidad en algún mes del año: Ruta 30.9\,\%, Bahía
24.2\,\%, Belice 38.3\,\%; concentrado en diciembre y enero.
\end{opciones}

\section*{El planeador de viaje (en la página)}
Brandon pidió: \emph{``como si alguien se mete a la pagina a planear sus vacaciones y el quiere ir a enero\dots{} y si ve
flujos super altos o temporada alta re recomiende otra parte\dots{} hospedagues hacer el NLP para recomendar donde,
hospedarse y restuaramntes\dots{} mediante google maps y webscraping NLP''}.

\begin{opciones}[5.9 · De dónde salen los lugares que se recomiendan · 1-oct-2026]
\begin{center}
\small
\begin{tabularx}{\linewidth}{>{\raggedright\arraybackslash}p{0.3\linewidth}Y}
\encabezado \h{Opción} & \h{Por qué sí o por qué no} \\
\textbf{DENUE + botón de Google Maps (elegida)} & Negocios oficiales con coordenadas; el botón es un enlace público que abre
Google Maps en otra pestaña. Sin scraping, sin llave, funciona sin internet \\
\filagris Scraping de Google Maps & Prohibido por sus términos y por la regla 4 \\
API de Google Places & Llave con tarjeta, deja de funcionar sin internet y la llave quedaría expuesta en un repositorio
público \\
\filagris Solo el mapa propio & 100\,\% sin internet, pero menos útil para llegar \\
\bottomrule
\end{tabularx}
\normalsize
\end{center}
\end{opciones}

\begin{opciones}[5.10 · Qué hace el NLP · 1-oct-2026]
\textbf{Elegido}: un clasificador por \textbf{léxico} sobre nombres y giros (tipo de lugar, grupo y momento del día).
\textbf{Por qué}: no hay etiquetas para entrenar un modelo estadístico; un léxico es honesto y se explica regla por regla.

\textbf{Resultado}: 1,587 de 2,042 negocios clasificados; \textbf{39 de 40} correctos en una muestra nueva revisada a mano.
\end{opciones}

\begin{opciones}[5.11 · Cuándo es temporada alta · 1-oct-2026]
\textbf{Elegida}: forma del año $\geq1.20$ \textbf{o} riesgo de rebasar la capacidad $\geq10\,\%$. \textbf{Descartado}: el
semáforo del Radar (es del mes actual, y los 5 lugares salen siempre tranquilos).

\textbf{Ajuste al probarlo}: sin una condición de lluvia, el mes sugerido casi siempre era junio, el más lluvioso (195 mm). Se
agregó que llueva menos que la mediana.
\end{opciones}

\begin{opciones}[5.12 · Cancún y Riviera Maya en el planeador · 1-oct-2026]
\textbf{Lo que pidió Brandon}: \emph{``mira esos dos como no tenemos datos cambialos por cancun y riviera maya''} (Maya Ka'an y
la Laguna, sin serie), y \emph{``Que aparezca, pero si Cancún está lleno o Riviera está lleno en ese mes que seleccionaron, le
recomendamos de manera que sí se vea y no pase desapercibido en un scroll que puede elegir otro lugar, y así no rompemos la
regla''}.

\textbf{Temporada alta del norte}: su ocupación cambia poco (Cancún, de 64.5\,\% a 80.7\,\%), así que con la regla del sur nunca
sería temporada alta (su índice máximo es 1.08). \textbf{Elegido}: los \textbf{cortes del Radar} (ocupación $\geq71.2\,\%$):
Cancún queda en temporada alta de noviembre a abril. \textbf{Descartados}: la regla del sur (nunca recomendaría el sur) y solo
``saturado'' $\geq85.9\,\%$ (ningún mes llega).

\textbf{Consecuencia}: la regla de oro 9 se ajustó. El norte aparece con la etiqueta ``Referencia: la campaña no lo
promueve'' y, si está lleno, el planeador \textbf{siempre} recomienda un lugar del sur. La Riviera Maya entró al Pronóstico con
el mismo método (Holt-Winters sin tendencia: el único con rango $\geq80\,\%$, aunque con 72\,\% más error que la línea base; se
declara).
\end{opciones}
